\documentclass[11pt]{letter}
\usepackage[margin=1in]{geometry}
\usepackage{hyperref}

\signature{Jeffrey D.\ Varner, Ph.D.\\
Professor\\
Robert Frederick Smith School of Chemical and Biomolecular Engineering\\
Cornell University\\
Ithaca, NY 14853\\
\texttt{jdv27@cornell.edu}}

\address{Robert Frederick Smith School of\\Chemical and Biomolecular Engineering\\Cornell University\\Ithaca, NY 14853}

\begin{document}

\begin{letter}{Editorial Office\\Frontiers in Bioengineering and Biotechnology}

\opening{Dear Editors,}

We are pleased to submit our manuscript entitled ``Mechanistic Resource Competition Modeling of a Cell-Free Gluconate Biosensor'' for consideration as an Original Research article in \textit{Frontiers in Bioengineering and Biotechnology}.

Cell-free synthetic biology systems are increasingly used for rapid prototyping of biosensor circuits, yet their finite transcriptional and translational resources impose time-dependent capacity constraints that are difficult to capture with empirical decay functions. In this study, we developed a mechanistic two-layer resource competition model for a D-gluconate biosensor circuit operating in the reconstituted cell-free system PURExpress. The first layer describes the allocation of RNA polymerase and ribosome between competing genes using a competitive enzyme kinetics formalism derived from elementary transcription and translation mechanisms. The second layer captures the irreversible depletion of consumable resources (nucleotides, energy cofactors, and amino acids) through ordinary differential equations driven by cumulative biosynthetic flux. This mechanistic formulation replaces the ad hoc logistic and exponential decay functions commonly used in cell-free models while providing a physical interpretation of resource exhaustion.

We believe this work makes three contributions that will be of interest to the readership of \textit{Frontiers in Bioengineering and Biotechnology}:

\begin{enumerate}
\item \textbf{A mechanistic resource framework for cell-free systems.} The two-layer model (machinery allocation + consumable depletion) provides a principled alternative to empirical decay functions, requiring only two resource consumption parameters ($\alpha_X$ and $\alpha_L$) compared to the multiple onset, slope, and half-life parameters needed by empirical approaches.

\item \textbf{Predictive capability without re-estimation.} The model accurately predicted the dose-response relationship of Venus protein as a function of D-gluconate concentration over four orders of magnitude, using only parameters estimated from the training conditions.

\item \textbf{Identification of the transcriptional bottleneck.} The model predicted that 76\% of transcriptional resources were consumed by 12 hours compared to less than 2\% of translational resources. A synthetic fed-batch supplementation analysis confirmed that replenishing transcriptional resources increased protein yield by up to 52\%, while replenishing translational resources had no detectable effect. These predictions were qualitatively consistent with independent experimental evidence from fed-batch studies in the PURE system.
\end{enumerate}

The model code, parameter ensemble, and experimental data are publicly available under an MIT license at \url{https://github.com/varnerlab/Gluconate-Sensor-Model-Paper}.

We suggest the following reviewers based on their expertise in cell-free systems modeling, resource competition theory, and biosensor design:

\begin{itemize}
\item Prof.\ Andras Gyorgy, New York University Abu Dhabi (\texttt{andras.gyorgy@nyu.edu}) --- resource competition theory for genetic circuits
\item Prof.\ Vincent Noireaux, University of Minnesota (\texttt{noireau@umn.edu}) --- cell-free gene expression systems and TX/TL dynamics
\item Prof.\ Sebastian J.\ Maerkl, EPFL (\texttt{sebastian.maerkl@epfl.ch}) --- PURE system characterization and performance limits
\item Prof.\ Domitilla Del Vecchio, Massachusetts Institute of Technology (\texttt{ddv@mit.edu}) --- resource-aware modeling of synthetic circuits
\item Prof.\ Enoch Yeung, University of California Santa Barbara (\texttt{eyeung@ucsb.edu}) --- mathematical modeling of cell-free TX/TL dynamics
\end{itemize}

This manuscript has not been published previously and is not under consideration elsewhere. All authors have approved the manuscript and agree with its submission to \textit{Frontiers in Bioengineering and Biotechnology}.

\longindentation=0pt
\closing{Sincerely,}

\end{letter}
\end{document}
